% chapter: 前提知識と計算の基本

\chapter{前提知識と計算の基本}\label{章：前提知識と計算の基本}\label{章：計算の基本}
\begin{bookcolumns}
  \bookterm{sequence}{数列}の読み方と、和や式変形の道具をまとめる。
  本編の計算で必要になったときに、対応する節へ戻れる。
  各節の冒頭には、その道具を使う場所を記した。

  \bookterm{logarithm}{対数}・\bookterm{trigonometric}{三角関数}・\bookterm{limit}{極限}・複素数の計算は、巻末の
  \chapterref{章：用語と計算の参照}にまとめた。
  本編でそれらが現れたところから、説明へ移れる。

\section{漸化式と数列の基本性質}
  まずは漸化式についてきちんと定義しておこう.
% section: シグマの計算

\section{シグマの計算}\label{節：シグマの計算}
  \noindent\textbf{【本編で使う場面】}\par
  和の記号と添字の役割を確認する.\sectionref*{節：階差型}で変化を足し戻し,\sectionref*{節：総和型}で総和の式をずらして引くときに使う.

  多くの項を足すとき、すべてを並べる代わりに、一つの記号でまとめられる。
  実数の\bookterm{sequence}{数列}$\{a_k\}$について、$a_1$から$a_n$までの和を
  \[
    \sum_{k=1}^{n} a_k=a_1+a_2+\cdots+ a_n
  \]
  と書く。複素数の\bookterm{sequence}{数列}でも、同じ記号を使う。
  $\sum$は、ギリシャ文字の\bookterm{sigma}{シグマ}である。
  下の$k=1$から上の$n$まで、$k$を一つずつ進めて足す。

  同じ数を繰り返し足せば、項数を掛けた値になる。
  定数倍は和の外へ出せる。二つの和は、項ごとに分けてもよい。
  \begin{align*}
    &\sum_{k=1}^{n} c=nc\\
    &\sum_{k=1}^{n} ca_k=c\sum_{k=1}^{n} a_k\\
    &\sum_{k=1}^{n} (a_k+b_k)=\sum_{k=1}^{n} a_k+\sum_{k=1}^{n} b_k
  \end{align*}
  自然数と、その二乗・三乗の和は、次の式でまとめられる。
  \begin{align*}
    &\sum_{k=1}^{n} k=\frac{1}{2}n(n+1)\\
    &\sum_{k=1}^{n} k^2=\frac{1}{6}n(n+1)(2n+1)\\
    &\sum_{k=1}^{n} k^3=\left\{\frac{1}{2}n(n+1)\right\}^2\\
  \end{align*}

  添字$k$は,和の中で順に動く番号である.同じ範囲なら$\sum_{k=1}^{n}a_k=\sum_{j=1}^{n}a_j$と書き換えても値は変わらない.
  一方,上端の$n$は足す項数を決める.例えば$S_n=\sum_{k=1}^{n}a_k$なら
  \[
    S_{n+1}=S_n+a_{n+1}
  \]
  である.
  \bookterm{recurrence}{漸化式}の添字をずらす場合は,式全体の$n$を置き換える.
  例えば$a_{n+1}=2a_n+n$から$a_{n+2}=2a_{n+1}+n+1$を得る.

% section: 式の変形

\section{式の変形}\label{節：式の変形}
  \noindent\textbf{【本編で使う場面】}\par
  同じ値を保ちながら既知の形を作る操作を確認する.定数の調整は\sectionref*{節：特性方程式型},\bookterm{partial}{部分分数分解}は\sectionref*{節：階差型}で和を計算する際に使う.

  \subsection{\texorpdfstring{$+a-a$をする}{+a-aをする}}\label{節：+a-aをする}
    二次式の中に、二乗の形を作る。
    $a\ne0$として、足りない定数を足し、同じだけ引くと、
    \begin{align*}
      &ax^2+bx+c\\
      &\qquad=a\left(x^2+\frac{b}{a}x\right) +c\\
      &\qquad=a\left(x^2+\frac{b}{a}x+\frac{b^2}{4a^2}-\frac{b^2}{4a^2}\right)+c \tag*{\(\cdots\) (☆)}\\
      &\qquad=a\left(x+\frac{b}{2a}\right)^2-\frac{b^2}{4a}+c
    \end{align*}
    となる。☆の行で、二乗に必要な定数を補った。
    足して引いた量は零なので、式の値は変わらない。
    $x$を含む部分が一つの二乗にまとまり、式の形が見やすくなる。
    この操作を\bookterm{square}{平方完成}という。
    二次関数の最大・最小や、軸の位置を調べるときにも使える。

    値を保ったまま、扱いやすい形を作る。
    同じ発想が、後の定数項を取り除く操作につながる。

  \subsection{部分分数分解}\label{節：部分分数分解}
    次は、積を含む分母を、より簡単な分数の和へ分ける。
    \[
      \frac{1}{(x+a)(x+b)}
    \]
    は、そのままでは和を計算しにくい。
    分母を一次の式に分けるため、次の形を作る。
    ここでは$x+a\ne0$、$x+b\ne0$とする。
    \[
      \frac{1}{(x+a)(x+b)}=\frac{\alpha}{x+a}+\frac{\beta}{x+b}
    \]
    $\alpha,\beta$について求めると,
    \begin{align*}
    \frac{1}{(x+a)(x+b)}
      &= \frac{\alpha}{x+a}+\frac{\beta}{x+b}\\
      &= \frac{(\alpha+\beta)x+\alpha b+\beta a}
      {(x+a)(x+b)}
    \end{align*}
    \[
    \therefore\quad
    \begin{cases}
    \alpha+\beta=0,\\
    \alpha b+\beta a=1
    \end{cases}
    \]
    これを解くと,\,$\beta=-\alpha$であるから,
    \[
      \alpha b-\alpha a=1
      \iff \alpha(b-a)=1
    \]
    したがって,\,$a\ne b$であれば,
    \[
      \alpha=\frac{1}{b-a},
      \qquad
      \beta=-\frac{1}{b-a}
    \]
    が得られる(\,$a=b$のときはそもそも分母が重解$(x+a)^2$になり,この形の分解はできない).
    二つの係数を定めると、一次の分母をもつ分数の和に分かれる。
    この操作を\bookterm{partial}{部分分数分解}という。
    3次の場合も,\,$p,q,r\ne0$で,分母の一次因子$px+a,qx+b,rx+c$の根
    \[
      -\frac{a}{p},\qquad -\frac{b}{q},\qquad -\frac{c}{r}
    \]
    が互いに異なるときは,
\BookMathLayout{\[
      \frac{1}{(px+a)(qx+b)(rx+c)}=\frac{\alpha}{px+a}+\frac{\beta}{qx+b}+\frac{\gamma}{rx+c}
    \]}{
\[
\begin{aligned}
\frac{1}{(px+a)(qx+b)(rx+c)}
&=\frac{\alpha}{px+a}+\frac{\beta}{qx+b}\\
&\quad+\frac{\gamma}{rx+c}
\end{aligned}
\]
}
    と置き,係数$\alpha,\beta,\gamma$を定めることで分解できる.
    言い換えれば,これらの一次因子が互いに定数倍でないことが条件である.
    重複因子があるときは,その重複度に応じて例えば$\dfrac{A}{x+a}+\dfrac{B}{(x+a)^2}$のような項も必要になる.

% section: 等差数列とその総和の導出

\section{等差数列とその総和の導出}\label{節：等差数列とその総和の導出}
  \noindent\textbf{【本編で使う場面】}\par
  等差数列の和を,並べ方の工夫から導く.\sectionref*{節：階差型}で一次式を足すときに使う.

  初項が$a$,公差が$d$である等差数列
  \[
    a,\,a+d,\,a+2d,\,\ldots,\,a+(n-1)d
  \]
  の初項から第$n$項までの総和を考える.
  その総和を$S_n$とし,第$n$項を$l$とすると,
  \[
    l=a+(n-1)d
  \]
  である.

  \subsection{総和の公式}\label{節：等差数列とその総和の導出：総和の公式}
    等差数列の初項から第$n$項までの総和は,
    \[
      S_n
      =
      \frac{n}{2}\{2a+(n-1)d\}
    \]
    である.
    第$n$項$l=a+(n-1)d$を用いれば,
    \[
      S_n=\frac{n}{2}(a+l)
    \]
    とも表せる.
    つまり,総和は項数と初項・末項の平均との積である.

  \subsection{公式の導出}\label{節：等差数列とその総和の導出：公式の導出}
  \BookFigArithmeticSum

    図は公差$d>0$の場合に,各項の増分
\BookMathLayout{$d,\,2d,\,\ldots,\,(n-1)d$と,それを逆順に並べたものを上下に重ねた模式図である.}{\\
$d,\,2d,\,\ldots$,$(n-1)d$と,それを逆順に並べたものを上下に重ねた模式図である.}
    上下を組み合わせた高さはどの列も$nd$になる.
    つまり,
    \[
      2\sum_{k=1}^{n-1}kd=n(n-1)d
    \]
    であり,共通の初項$a$を$n$回加えれば,
    $S_n=na+\displaystyle\frac{n(n-1)}{2}d$となる.
    以下では,同じ工夫を元の等差数列そのものに適用する.


    総和を順に並べた式と,逆順に並べた式を書く.
    \[
      \begin{array}{rcl}
        S_n
          &=&
          a+(a+d)+(a+2d)\\
          &&\qquad +\cdots+\{a+(n-1)d\}\\
        S_n
          &=&
          \{a+(n-1)d\}+\{a+(n-2)d\}\\
          &&\qquad +\cdots+(a+d)+a
      \end{array}
    \]
    これらを項ごとに足すと,各列の和は$2a+(n-1)d$となる.
    そのような列が$n$個あるので,
    \[
      2S_n=n\{2a+(n-1)d\}
    \]
    である.
    両辺を$2$で割れば,
    \[
      S_n=\frac{n}{2}\{2a+(n-1)d\}
    \]
    を得る.

    この導出では,同じ数列を逆順に並べて加えることで,各列の和を一定にしている.


% section: 等比数列の総和とその導出

\section{等比数列の総和とその導出}\label{節：等比数列の総和とその導出}
  \noindent\textbf{【本編で使う場面】}\par
  等比数列の和を,倍率をそろえて引く操作から導く.\sectionref*{節：階差型}と\sectionref*{節：特性方程式型}の階差による別解で使う.

  初項が$a$,公比が$r$である等比数列
  \[
    a,\,ar,\,ar^2,\,\ldots,\,ar^{n-1}
  \]
  の初項から第$n$項までの総和を考える.
  その総和を$S_n$とすると,
  \[
    S_n
    =\sum_{k=1}^{n} ar^{k-1}
    =a+ar+ar^2+\cdots+ar^{n-1}
  \]
  と表せる.

  \subsection{総和の公式}\label{節：等比数列の総和とその導出：総和の公式}
    公比$r$が$1$でないとき,
    \[
      S_n=\frac{a(1-r^n)}{1-r}
    \]
    である.
    分母と分子の符号を同時に変えれば,
    \[
      S_n=\frac{a(r^n-1)}{r-1}
    \]
    とも書ける.
    公比が$1$のときは,すべての項が$a$であるから,
    \[
      S_n=na
    \]
    となる.

  \subsection{公式の導出}\label{節：等比数列の総和とその導出：公式の導出}
    $r\ne1$として,総和の式の両辺に$r$を掛ける.
    \begin{align*}
      S_n
        &=a+ar+ar^2+\cdots+ar^{n-1},\\
      rS_n
        &=ar+ar^2+ar^3+\cdots+ar^n.
    \end{align*}
    2つの式を項ごとに引くと,右辺では共通する項がすべて消える.
\BookMathLayout{\[
    \begin{array}{rrrl}
        &S_{n}={}&a+&\hspace{-0.75em}ar+ar^2+\cdots+ar^{n-1}\\
      -\,) & rS_n={}&{}&\hspace{-0.75em}ar+ar^2+ar^3+\cdots+ar^{n}\\ \hline
        & (1-r)S_n={}&{}&\hspace{-0.75em}a-ar^n
    \end{array}
    \]}{
\[
\begin{array}{rl}
S_n={}&a+ar+ar^2+\cdots\\
     &{}+ar^{n-1}\\
-\,)\ rS_n={}&ar+ar^2+ar^3+\cdots\\
     &{}+ar^n\\ \hline
(1-r)S_n={}&a-ar^n
\end{array}
\]
}
    よって,
    \[
      (1-r)S_n=a(1-r^n)
    \]
    である.
    $r\ne1$より両辺を$1-r$で割ると,
    \[
      S_n=\frac{a(1-r^n)}{1-r}
    \]
    を得る.

    この導出では,総和を公比$r$倍してから元の総和との差を取ることで,途中の項を消している.
    数列の和を求めるときには,このように総和を定数倍して引き,共通する項を消す考え方がよく用いられる.


% section: 数学的帰納法

  \section{数学的帰納法}\label{節：数学的帰納法}
  \noindent\textbf{【本編で使う場面】}\par
  数項からの予想をすべての項の証明へ進める方法を学ぶ.\sectionref*{節：漸化式を解く前に}の検算や,\sectionref*{節：虚数解型隣接3項間漸化式}の\bookterm{general}{一般項}の確認で使う.

    ここまでは計算に用いる公式や式変形を扱った. ここでは,自然数に関する主張を証明する方法を説明する.
    自然数$n$についての主張$P(n)$が,すべての$n\ge1$で成り立つことを示したいとする.
    このとき,以下の2つを示せばよい.
    \begin{enumerate}[label=(\roman*),parsep=3pt]
      \item $P(1)$が成り立つ.
      \item $n\ge1$を任意にとり,\,$P(n)$が成り立つと仮定すると,\,$P(n+1)$も成り立つ.
    \end{enumerate}
    この2つが示されれば,\,(i)より$P(1)$が成り立ち,\,(ii)より$P(1)\Rightarrow P(2)$,\,$P(2)\Rightarrow P(3)$,\,$\ldots$と次々に成り立つことが分かるので,すべての$n\ge1$について$P(n)$が成り立つ.
    ドミノ倒しに例えられることが多い.\,(i)は最初のドミノを倒すこと,\,(ii)は「どのドミノについても,それが倒れれば次のドミノも倒れる」ように並べることに対応する.

    \sectionref*{節：極形式と複素平面}では,この方法を用いて\bookterm{de-moivre}{ド・モアブルの定理}を証明する.
    また,\,\chapterref{章：漸化式の解説}では,\bookterm{recurrence}{漸化式}から予想した\bookterm{general}{一般項}の形を確かめる場面や,\bookterm{trigonometric}{三角関数}の\bookterm{addition}{加法定理}から\bookterm{recurrence}{漸化式}を導く場面(\sectionref{節：漸化式を解く前に}や\sectionref{節：虚数解型隣接3項間漸化式}など)で\bookterm{induction}{数学的帰納法}を用いる.

\end{bookcolumns}
